
\section*{2. Premières réductions}

\textbf{2.1.} Appliquons 1.3 à @@M0@@ et @@M1@@~: on trouve un ouvert non vide @@M2@@ de @@M3@@, un schéma intègre @@M4@@ et un morphisme fini plat @@M5@@ tels que toute composante irréductible @@M6@@ de @@M7@@ et @@M8@@ de @@M9@@, munie de la structure réduite induite, soit intègre sur @@M10@@. En vertu de la descente f.p.q.c. (EGA IV 2.7.1 (xiv) et SGA 1 VIII 7.9), on peut remplacer @@M11@@ par @@M12@@. Il résulte alors du diagramme commutatif suivant~:
@@M13@@
impliquant le diagramme commutatif
@@M14@@
et de sorites triviaux sur les morphismes de foncteurs représentables par morphismes quasi-affines (résultant de EGA II 5.1.10 (ii) (v)) qu'il suffit de démontrer deux assertions suivantes :
\begin{enumerate}[label=\Roman*.]
\item \uline{Le théorème} 1.1 \uline{lorsque} @@M15@@ \uline{est une immersion (nécessairement nilpotente).}
\item \uline{Le théorème} 1.1 \uline{lorsque} @@M16@@ \uline{et} @@M17@@ \uline{sont réduits et que toute composante irréductible} @@M18@@ \uline{de} @@M19@@ \uline{et} @@M20@@ \uline{de} @@M21@@\uline{, munie de la structure réduite induite, est intègre sur} @@M22@@.
\end{enumerate}

\textbf{2.2.} Étudions l'assertion II. Pour tout @@M23@@, soit @@M24@@ un indice tel que @@M25@@, soit @@M26@@. Alors on a les deux diagrammes
